\section{Lecciones para la cola de trabajo siguiente: por qué EP139 debe servir de modelo para los episodios de revisión}

Este episodio es de tipo review/survey, no una entrevista a una persona. Ofrece a la cola de trabajo siguiente de Zhang Xiaojun AI una plantilla de procesamiento: cuando hay slides, se hace slide-complete; cuando no hay slides, no se insertan una y otra vez fotogramas de las personas, sino que se generan diagramas conceptuales, glosarios y diagramas de mecanismos. En los programas de revisión, el valor de la nota está en reconstruir la estructura del conocimiento, no en condensar la conversación.

\begin{knowledgebox}{Plantilla de generación para pódcasts de revisión}
Primero, extraer la genealogía técnica; segundo, identificar los términos centrales; tercero, convertir los juicios expresados de forma oral en tablas o diagramas de flujo; cuarto, señalar las controversias y los límites; quinto, usar la sección de síntesis para enlazar con el siguiente episodio o el siguiente tema. Solo así la nota resultante se parece a un material de clase y no a una transcripción editada.
\end{knowledgebox}

\subsection{Resumen de la sección}

La metodología de EP139 se aplicará a la cola de trabajo siguiente sobre IA e internet en sentido amplio: en el cuerpo del texto solo entran la portada y las imágenes que de verdad aportan información; cuando no hay contenido visual, la estructura didáctica se apoya en gráficos generados; en cada episodio se conservan transcript, manifest, coverage y visual QA.


